\chapter*{Introduction}
\addcontentsline{toc}{chapter}{Introduction}

\section*{Relevance of the Research}

The transition of electric power systems toward low-carbon operation has turned power system simulation and electricity market optimization into a computational bottleneck rather than a modeling formality. The large-scale integration of variable renewable generation, energy storage, and sector coupling has multiplied the number of decision variables, time steps, and scenarios that a single study must evaluate. Planning a decarbonization pathway for a national grid now routinely requires solving optimization problems with millions of variables across thousands of representative hours, and repeating that solve for dozens of policy or weather scenarios.

In the academic and open-research community, PyPSA (Python for Power System Analysis) has become the de facto standard framework for this class of problems. It underpins flagship models such as PyPSA-Eur and PyPSA-Earth and supports a large body of peer-reviewed work on renewable energy systems. Its strength is a clean, high-level modeling layer built on the Python scientific stack. That same layer, however, is also the source of a practical limitation: for large networks and iterative scenario studies, the time spent building the optimization model in Python --- assembling constraints through pandas and the algebraic-modeling layer --- can dominate the total runtime and reach hours per solve, well before the numerical solver itself begins work.

The Julia programming language offers a direct way to address this layer. Julia combines a high-level, dynamically typed syntax with just-in-time compilation to native code, and its JuMP modeling framework compiles the constraint-assembly step instead of interpreting it. This makes Julia a credible candidate for reimplementing the modeling layer that limits PyPSA, while delegating the actual numerical solving to the same industrial solvers (HiGHS, Ipopt) that PyPSA already uses.

At the same time, the operational use of these models increasingly depends on forecasts of uncertain quantities --- most importantly electricity demand. A dispatch schedule that is optimal for a single point forecast can be far from optimal, or even infeasible, once the realized demand deviates from that forecast. This motivates coupling data-driven forecasting with the physics-based optimization, so that the optimizer is informed not only by a predicted demand profile but also by a calibrated description of its uncertainty. The combination of a high-performance Julia solver layer with an uncertainty-aware machine-learning forecast defines the ``AI-assisted'' framing of this work and explains its relevance to both the power systems and the scientific computing communities.

\section*{Problem Statement}

Despite Julia's well-documented performance characteristics, no actively maintained, self-contained reimplementation of PyPSA's core computational modules currently exists, and the question of how much of PyPSA's runtime is actually recoverable by changing the modeling layer has not been answered with reproducible measurements. Equally, the integration of a statistically validated demand forecast into a stochastic dispatch formulation within a single native framework remains largely unexplored. The present work is organized around the following concrete questions:

\begin{itemize}
    \item Can PyPSA's core algorithms --- DC and AC power flow, linear optimal power flow, and unit commitment --- be reimplemented in Julia while preserving numerical equivalence with the Python reference?
    \item What speedup is actually attributable to the Julia modeling layer, once both implementations are made to call the same underlying solver, and how does that speedup scale with network size and planning horizon?
    \item Can a neural load forecaster be equipped with distribution-free prediction intervals that provide a guaranteed coverage level suitable for dispatch optimization?
    \item Does propagating forecast uncertainty into a scenario-based stochastic dispatch yield decision-relevant risk measures that a deterministic formulation cannot provide?
\end{itemize}

\section*{Object and Subject of Research}

\textbf{The object of research} is the set of computational methods used for steady-state power system simulation and for the optimization of electricity market dispatch, together with the data-driven forecasting methods that supply their demand inputs.

\textbf{The subject of research} is the implementation of these methods in the Julia programming language with a PyPSA-compatible interface, the quantitative comparison of their performance against the PyPSA reference implementation, and the integration of an uncertainty-aware load forecast into a stochastic dispatch formulation within that implementation.

\section*{Purpose of the Research}

The purpose of this thesis is to design, implement, and validate a self-contained Julia framework that reproduces the core simulation and optimization capabilities of PyPSA with measurable performance gains, and to extend it with an AI-assisted layer that forecasts electricity demand under uncertainty and propagates that uncertainty into the dispatch optimization.

\section*{Research Tasks}

To achieve this purpose, the following tasks are addressed:

\begin{enumerate}
    \item Implement and validate DC and AC power flow solvers in Julia, establishing numerical equivalence with PyPSA on standard test networks.
    \item Implement single-period and multi-period Linear Optimal Power Flow (LOPF) with energy storage and variable renewable generation, including the extraction of locational marginal prices from the dual solution.
    \item Extend the LOPF to a Unit Commitment formulation as a mixed-integer linear program with minimum up/down times and start-up costs.
    \item Design a reproducible benchmarking methodology and quantify the speedup of the Julia implementation over PyPSA across all solver types, a range of network sizes, and several planning horizons.
    \item Develop a machine-learning load forecaster, equip it with conformal prediction intervals that guarantee a prescribed marginal coverage, and integrate the resulting demand scenarios into a sample-average-approximation stochastic LOPF.
\end{enumerate}

\section*{Scientific Novelty}

The scientific novelty of the work consists in the following:

\begin{itemize}
    \item A self-contained, validated Julia reimplementation of PyPSA's core computational modules with a PyPSA-compatible construction API, requiring no Python dependency at solve time, in contrast to prior partial ports that retained Python in the workflow.
    \item A reproducible, multi-seed decomposition of the Julia-versus-PyPSA performance difference that attributes the speedup to the model-assembly layer rather than to the numerical solver, reported as compute-bound figures with means and standard deviations rather than headline ratios.
    \item The integration, within a single native framework, of a conformal-prediction-calibrated demand forecast with a scenario-based stochastic dispatch, producing distribution-free coverage guarantees and decision-relevant risk measures (expected cost and conditional value-at-risk) that are unavailable to deterministic dispatch.
\end{itemize}

\section*{Practical Significance}

The resulting \texttt{PowerFlowJulia} framework provides a drop-in path for PyPSA users to accelerate the model-building stage of large or iterative studies without changing their numerical backend, and demonstrates a concrete pipeline in which a calibrated forecast informs risk-aware dispatch decisions. The implementation, benchmark scripts, and data pipeline are released as a reproducible package.

\section*{Thesis Structure}

The thesis follows a survey--design--implement--evaluate arc. Chapter~\ref{ch:litreview} reviews the state of the art: the PyPSA framework and its modeling layer, the existing Julia power-system ecosystem, and the relevant literature on load forecasting and stochastic dispatch, and it identifies the research gap addressed here. Chapter~\ref{ch:methodology} develops the mathematical and methodological foundations of every implemented method --- the DC and AC power flow models, the LOPF and unit commitment formulations, the load-forecasting model with conformal prediction, and the stochastic LOPF --- together with the benchmarking protocol. Chapter~\ref{ch:implementation} presents the Julia implementation of each module, validates it against PyPSA, and reports the performance benchmarks and the results of the AI-assisted dispatch experiments. Chapter~\ref{ch:conclusions} synthesizes the findings, states the conclusions regarding Julia's suitability as a high-performance substitute for PyPSA and the value of AI-assisted dispatch, and outlines directions for future work.
